% ============================================================
%  Lecture 1: The stage (curve, group, moduli space)
% ============================================================

\section{舞台：一条曲线、一个群、一个模空间}
\label{sec:stage}

\begin{goals}
\begin{enumerate}
\item 说清楚这十讲的主角为什么是一个三元组 $(C,G,\MH)$，其中每个字母各管什么。
\item 会写 Hitchin 方程与 Hitchin 对 $(E,\phi)$ 的定义，并能解释 $\phi$ 为什么既是“场”又是“一次形式”。
\item 能说清 $\MH$ 的三个复结构、维数公式，以及 Hitchin 纤维化为什么让二维 sigma 模型变得可解。
\end{enumerate}
\end{goals}

\subsection{为什么是一条曲线和一个群}

先把结论放在这里：后面十讲反复出现的那个对象，是一个二维 sigma 模型的 target，而它由两样东西决定——一条紧 Riemann 面 $C$，和一个紧 Lie 群 $G$。这两样东西分别来自几何与物理，本讲其余部分都在解释它们怎么合到一起。

几何那头，主角是 $C$ 上的 $G$-丛，等价地是 $\pi_1(C)\to G$ 的表示。几何 Langlands 纲领说的是：$C$ 上 $G$ 侧的对象与 $\Ldual$ 侧的对象之间有一一对应。这个陈述本身是纯几何的，与物理无关；\kw{} 做的事情（\cite{kw06}）是给它一个物理来源。

物理那头，四维 $\mathcal N=4$ 超对称 Yang--Mills 理论先做一次特殊的拓扑扭化（第 3 讲详述），再放到乘积 $M=\Sigma\times C$ 上。从 $\Sigma$ 的方向看，低能理论是一个二维 sigma 模型，target 就是 $\MH$。原文的说法是：四维扭化理论 on $\Sigma$ reduces to a sigma-model of maps $\Phi:\Sigma\to \MH$ \src{B p.38}。这句话是整份讲义的起点，你把它记住，理由后面都会补上。

为什么 $C$ 要紧、而且要 $g\ge 2$？有两个理由。第一，$g\ge2$ 时半单 $G$ 上的平坦联络一般不可约，$\MH$ 的一般点也就光滑，可约解给出的奇异点可以被躲开 \src{B p.39}。第二，$C$ 紧才使后面要用的紧性成立：Hitchin 纤维化的纤维是紧的复环面，“纤维紧”这一条在原文里反复出现。反过来，若 $C$ 是球面，$\pi_1$ 平凡；若是环面，$2g-2=0$，维数公式直接退化。所以 $g\ge2$ 不是技术性的偷懒，而是舞台本身的要求。

$\mathcal N=4$ 这个“四”也不是随意的：它给出十六个超荷，正好是四维无引力规范理论能带的最大超对称（第 2 讲细说）。而 $G$ 要半单、要紧：四维规范理论要酉性，规范群就得紧；几何描述里到处用复化 $G_\CC$，所以原文一开头就声明 always assume $G$ to be compact and denote its complexification as $G_\CC$ \src{B p.6}。至于半单——若 $G$ 含 $U(1)$ 因子，每个解都带连续的 $U(1)$ 规范对称，低能描述会变成 sigma 模型乘一个超对称 $U(1)$ 规范理论，原文明确说自己只在半单情形工作 \src{B p.40}。这两条约定你先收下，本讲不再提。

\subsection{丛、Higgs 场，以及 $\phi$ 的两个身份}

定义从一个 $G$-丛 $E\to C$ 开始，$A$ 是它的联络。这里有一个二维特有的魔术：在复维 $1$ 上，$A$ 的 $(0,1)$ 部分 $\bar\partial_A$ 自动满足 $\bar\partial_A^2=0$——因为 $(0,2)$-形式在曲线上不存在。于是任意联络都把 $E$ 变成全纯丛。原文一句话带过：any connection on a smooth $G$-bundle over a Riemann surface automatically turns $E$ into a holomorphic $G$-bundle \src{B p.49}。高维没有这种好事，所以“联络即全纯结构”是这一讲的第一个礼物，也是为什么整套构造只在曲线上成立。

第二个对象是 \term{Higgs field}{Higgs 场} $\phi$：它是 $K_C\otimes\mathrm{ad}(E)$ 的全纯截面，也就是取值在伴随丛里、带一个 $(1,0)$ 一次形式的全纯对象 \src{B p.49}。把 $E$ 连同 $\phi$ 一起叫做 \term{Higgs bundle}{Higgs 丛}。$g\ge2$ 时 $H^0(C,K_C\otimes\mathrm{ad}E)$ 非零，所以这不是空定义。

现在说本讲第一个要点：$\phi$ 有两个身份，缺一个都读不下去。

\textbf{身份一，它是场。}四维理论里有六个实标量，扭化之后其中四个被重新解释为四维流形上的一次形式；紧化到 $C$ 上，它的 $(1,0)$ 部分就是 $\phi$ \src{B p.56}。所以 $\phi$ 不是标量，也不是凭空造出来的几何对象——它来自四维场论。这个身份在第 3 讲扭化时会再出现一次，那时你会知道“为什么恰好是四个标量、为什么恰好是一次形式”。

\textbf{身份二，它是余切方向。}稳定 $G$-丛的模空间记作 $M$（原文的 $\mathcal M$），它在 $E$ 处的切空间是 $H^1(C,\mathrm{ad}E)$；由 Serre 对偶，余切空间是 $H^0(C,K_C\otimes\mathrm{ad}E)$。而 $\phi$ 正好住在这里。于是 $\phi$ 是“$M$ 的一个余切方向”：$\phi=0$ 的切片就是 $M\subset\MH$，更一般地 $\MH$ 包含 $T^*M$ 作为开稠密子集（差一个高余维的坏集）\src{B p.52}。

两个身份合起来，图像是“一个丛，加一个余切方向”。这也预告了下一节：三个复结构里，总有一个把 $\phi$ 当几何数据（Higgs 丛），另一个把 $\phi$ 与 $A$ 混成复联络（平坦联络）。

这里容易卡住：不要把 $\phi$ 当标量。一旦当标量，复结构 $J$ 里“$A+\ii\phi$ 是复联络”的写法立刻失去意义——$A$ 是一次形式，$\phi$ 必须也是一次形式，两者才能相加 \src{B p.51}。

\subsection{三个复结构，以及为什么存在两种商}

把所有的 $(A,\phi)$ 收进一个无穷维仿射空间 $W$。它带一个平坦度规（原文 (4.2)，只需知道它把 $\Tr(\delta A\wedge*\delta A+\delta\phi\wedge*\delta\phi)$ 在 $C$ 上积分）\src{B p.41}。$W$ 上可以定义三个复结构 $I,J,K$；用“哪些线性函数全纯”来定义，比写矩阵清楚：

\begin{itemize}
\item $I$：全纯坐标是 $A^{0,1}$ 与 $\phi^{1,0}$，即 $A$ 的 $(0,1)$ 部分与 $\phi$ 的 $(1,0)$ 部分全纯；
\item $J$：全纯坐标是 $A+\ii\phi$，即 $A$ 与 $\ii\phi$ 的各分量一起；
\item $K$：全纯坐标是 $A_z-\phi_z$ 与 $A_{\bar z}+\phi_{\bar z}$。
\end{itemize}

这三个结构由四元数关系 $IJ=-JI=K$ 等联系在一起。原文把它们推广成一整族复结构 $I_w$，参数 $w\in\CC\PP^1_h$：$w=0$ 是 $I$，$w=\infty$ 是 $-I$，$w=-\ii$ 是 $J$，$w=1$ 是 $K$ \src{B p.43 eq 4.6}。另有一个 $U(1)$ 作用 $\phi\mapsto\lambda\phi$（$|\lambda|=1$）把 $I_w$ 送到 $I_{\lambda^{-1}w}$，并固定 $w=0,\infty$ 两点 \src{B p.43}。请记住：$w$ 只是“$\MH$ 上复结构族”的参数，不是理论参数；理论参数 $\Psi$ 要到第 3 讲才出现，把这两个混起来是最常见的误读。

三个复结构各自配一个辛形式 $\omega_I,\omega_J,\omega_K$，以及 $\Omega_I=\omega_J+\ii\omega_K$ 等 \src{B p.44 eq 4.8}。这里有一个不对称，值得停下来看一眼：$\omega_I$ 与 $\omega_K$ 可以用微分形式整体写出（不含 $C$ 的复结构），而 $\omega_J$ 不行；对 $\Omega$ 来说关系正好反过来——$\Omega_I,\Omega_K$ 依赖 $C$ 的复结构，$\Omega_J$ 不依赖。这个不对称在讲镜像与 S-对偶时是主角，现在只要记住它存在。

现在做商。规范群作用在 $W$ 上，保持这三个辛形式；它给出的三个矩映射，其零点正好是 Hitchin 方程的三条：
\[
\mu_I=\mu_J=\mu_K=0
\quad\Longleftrightarrow\quad
F-\phi\wedge\phi=0,\qquad D\phi=D^*\phi=0 ,
\]
这两边分别是原文的 (4.11) 与 (4.1) \src{B p.46, B p.39}。于是
\[
\MH=W/\!/\!/\,G
\]
是一个 \term{hyper-Kähler quotient}{超凯勒商} \src{B p.46}。这个名词是字面意思：三个矩映射同时为零再商，商出来的流形自动带三个复结构——所以“超凯勒”不是额外假设，而是商法的产物。

这里有一个值得花五分钟想的区别：为什么不能只做一个复商？原文把两种商摆在一起看 \src{B p.49}。设 $\nu_I=\mu_J+\ii\mu_K$ 是复结构 $I$ 下的全纯矩映射，$\nu_I=0$ 就是 $\bar\partial_A\phi=0$（即 $\phi$ 全纯）。在它的零点集上有两个选择：

\begin{enumerate}
\item[(I)] 再要求实矩映射 $\mu_I=0$，只商 $G$-值的规范变换；
\item[(II)] 直接商 $G_\CC$-值的规范变换。
\end{enumerate}

(I) 给出带三个复结构的超凯勒流形；(II) 给出的只是复流形，而且只在当前这个复结构里有意义。两者相差一个高余维的稳定化修正，但在“看哪个复结构”这件事上完全不同：(II) 在 $I$ 下把对象变成稳定 Hitchin 对 $(E,\phi)$，这是 Hitchin 的定理 \src{B p.50}；(II) 在 $J$ 下把对象变成平坦 $G_\CC$ 联络，亦即 $\pi_1(C)\to G_\CC$ 的稳定同态，这是 Corlette--Donaldson 的定理 \src{B p.51}。而所有 $w\ne0,\infty$ 的 $I_w$ 给出的是同一个空间 $Y$ \src{B p.52}。所以“同一个模空间，三种复结构”不是比喻，而是两条定理拼起来的事实。表 \ref{tab:three-I} 把它压成一张表。

\begin{table}[htbp]
\centering
\small
\begin{tabular}{@{}llll@{}}
\toprule
复结构 & 全纯坐标 & 矩映射为零 & 商 $G_\CC$ 后的对象 \\
\midrule
$I$（$w=0$） & $A^{0,1},\ \phi^{1,0}$ & $\bar\partial_A\phi=0$ & 稳定 Hitchin 对 $(E,\phi)$ \\
$J$（$w=-\ii$） & $A+\ii\phi$ & $F(A+\ii\phi)=0$ & 稳定平坦 $G_\CC$ 联络 \\
$K$（$w=1$） & $A_z-\phi_z,\ A_{\bar z}+\phi_{\bar z}$ & 复联络 $(A_z-\phi_z,A_{\bar z}+\phi_{\bar z})$ 平坦 & 同一个 $Y$ \\
\bottomrule
\end{tabular}
\caption{$\MH$ 的三个复结构：同一个点集，三种全纯描述。前两行分别是 Hitchin 定理与 Corlette--Donaldson 定理，出处见 \src{B pp.50--52}。}
\label{tab:three-I}
\end{table}

\subsection{维数、Hitchin 纤维化，以及“为什么可解”}

先记结论：$\MH$ 是复维 $(2g-2)\dim G$ 的复流形 \src{B p.39}。本讲不证。直觉是：每多一个把手，就多 $2\dim G$ 个复方向；$g=2$、$G=SU(2)$ 时是 $6$，而稳定丛模空间 $M$ 是 $3g-3=3$ 维，$T^*M$ 正好 $6$ 维，与上一节“$\MH\supset T^*M$”对得上。原文也只是引用 Hitchin 的结果，不在这里重算。

真正让人兴奋的是 \term{Hitchin fibration}{Hitchin 纤维化}。取 $\phi$ 的 $G$-不变多项式 $P_i$，$i=1,\dots,r$（$r=\mathrm{rank}\,G$），次数 $d_i$ 满足
\[
\sum_{i=1}^{r}(2d_i-1)=\dim G .
\]
这些多项式给出映射
\[
\pi:\ \MH\longrightarrow B:=\bigoplus_{i=1}^{r}H^0\bigl(C,K_C^{\,d_i}\bigr),
\qquad (E,\phi)\longmapsto\bigl(P_1(\phi),\dots,P_r(\phi)\bigr),
\]
公式 (4.24) 见 \src{B p.53}。$SU(2)$ 时 $B=H^0(C,K_C^2)\cong\CC^{3g-3}$，映射就是 $\phi\mapsto\Tr\phi^2$ \src{B p.53}。因为 $\dim H^0(C,K_C^{\,d})=(2d-1)(g-1)$，从上面的恒等式直接算出
\[
\dim B=(g-1)\dim G=\tfrac12\dim \MH ,
\]
所以一般纤维的复维数也是 $(g-1)\dim G$ \src{B p.54 eq (4.25)}。原文还给了一个漂亮的计数：$SU(2)$ 时对固定 $w\in B$，方程 $\Tr\phi^2=w$ 是 $3g-3$ 个二次方程对 $3g-3$ 个未知数，一般有 $2^{3g-3}$ 个解 \src{B p.54}。

接下来是 Hitchin 的两条主定理。第一条：$\MH$ 在复结构 $I$ 下是 \term{completely integrable system}{完全可积系统}——存在 $\tfrac12\dim\MH$ 个 $I$-全纯、代数独立、对 $\Omega_I$ 的 Poisson 括号两两交换的函数。它们可以取成 $B$ 上线性函数的拉回，显式写法是
\[
H_{\alpha,i}=\int_C \alpha\wedge P_i(\phi),\qquad \alpha\in H^1\bigl(C,K_C^{\,1-d_i}\bigr),
\]
见 \src{B p.55 eq (4.26)}。交换性有一个一句话证明：在 $W$ 上 $\phi_z$ 只与 $A_{\bar z}$ 共轭、与 $\phi_z$ 自己交换，而 $H_{\alpha,i}$ 只是 $\phi$ 的函数 \src{B p.55}。第二条：这些对易流给出 $I$-全纯的交换流，而复环面恰好承认交换流；再由全纯辛几何的一般结果与纤维紧性，得到一般纤维是复环面 \src{B p.55}。原文还补了一句本单元最值钱的结论：纤维由 $H_k=h_k$ 定义，维数是一半、并且 $\Omega_I$ 在其上各向同性，所以纤维是 Lagrange 子流形 \src{B p.56}。

顺带说一句：$(2g-2)\dim G$ 与 $(g-1)\dim G$ 这两个数会反复出现，因为 $\MH$ 的“一半”在底 $B$ 上，“一半”在纤维里。以后看到“一半维数”的说法，回到这一页。

最后回答本讲标题里的问题：这件事为什么让二维 sigma 模型可解？当一个 sigma 模型的 target 是“底加紧环面纤维”的纤维化时，在纤维方向做 T-对偶几乎不改变物理，只是把周期取倒数。于是同一个二维理论可以在一个复结构里看成 $B$ 模型、在另一个里看成 $A$ 模型，两者互为镜像；而 $G$ 与 $\Ldual$ 的 Hitchin 系统恰好互为镜像 \cite{w08}。这正是原文第 5 讲“四维电磁对偶降成二维 T-对偶”的机制。

但这里必须诚实，这是本讲最容易过度宣称的地方：原文 §4.3 只造好了 $G$ 这一侧——纤维是复环面、纤维 Lagrange、纤维化存在，这些都是本节原话；而 T-对偶、镜像这些词在 §4 里根本没出现，它们要到 §5.3 才第一次出场 \src{B p.70}。所以准确的说法是：§4 给的是“可解的结构”，而不是“已经解了”。你读原文 §4 时，不要提前把它当成推导已完成。

\begin{figure}[htbp]
\centering
\begin{tikzpicture}[>=Stealth, font=\small]
  % 底 B
  \draw[kwblue,thick] (-4.0,0) -- (4.0,0);
  \node[kwblue,anchor=west] at (4.05,0) {$B$};
  % 一般纤维：环面画成带孔的椭圆
  \foreach \x in {-2.9,-1.3}
  {
    \draw[kwred,thick] (\x,1.85) ellipse (0.6 and 0.4);
    \draw[kwred,thick] (\x,1.85) ellipse (0.2 and 0.13);
    \draw[->,kwgray] (\x,0.12) -- (\x,1.35);
    \fill[kwblue] (\x,0) circle (1.2pt);
  }
  \node[font=\scriptsize,kwred] at (-2.1,2.55) {一般纤维：复环面（Lagrange）};
  % 退化纤维：8 字形
  \draw[kwred,thick] (1.7,2.2) circle (0.34);
  \draw[kwred,thick] (1.7,1.52) circle (0.34);
  \draw[->,kwgray] (1.7,0.12) -- (1.7,1.15);
  \draw[kwred,fill=white] (1.7,0) circle (2pt);
  \node[font=\scriptsize,kwred] at (1.7,2.75) {判别式轨迹};
  % 维数标注
  \node[font=\footnotesize,anchor=west] at (-3.95,-0.55) {底 $B$：复维 $(g-1)\dim G$};
  \node[font=\footnotesize,anchor=west] at (0.75,-0.55) {纤维：复维 $(g-1)\dim G$};
  \node[left=2pt,font=\scriptsize] at (-1.28,0.7) {$\pi$};
\end{tikzpicture}
\caption{Hitchin 纤维化 $\pi:\MH\to B$。一般纤维是复环面，并且在 $\Omega_I$（等价地 $\omega_J,\omega_K$）下是 Lagrange 子流形；在判别式轨迹上纤维退化。底与纤维各占一半复维数，这是后文“可解”的全部几何来源。}
\label{fig:fibration}
\end{figure}

\begin{keybox}
$\MH=W/\!/\!/\,G$ 是无穷维仿射空间按三个矩映射（即 Hitchin 方程）做的超凯勒商；它同时带三个复结构、复维 $(2g-2)\dim G$，并被 Hitchin 纤维化成底为 $B$（复维一半）、纤维为复环面（Lagrange）的可积系统。二维 sigma 模型的可解性就建立在这“一半环面加一半对易哈密顿量”的几何上——而对偶那一半要到 §5 才补上。
\end{keybox}

\begin{follow}
\begin{itemize}
\item \textbf{B §4 开头两页，pp.38--39}：紧化、极小解、Hitchin 方程 (4.1) 与维数公式。这是全讲的出发点，值得先读三遍。
\item \textbf{B §4.1，pp.41--44}：$W$ 的度规 (4.2)、$I,J,K$ 的矩阵定义 (4.3)--(4.5)、$I_w$ 家族 (4.6)。不必记矩阵，记住“哪个量依赖 $C$ 的复结构”。
\item \textbf{B §4.2，pp.49--52}：两种商 (I)/(II) 与两条定理（Hitchin、Corlette--Donaldson）。本讲第 3 节的内容基本都在这一段里。
\item \textbf{B §4.3，pp.52--56}：纤维化 (4.24)--(4.25)、完全可积 (4.26)、纤维是 Lagrange。p.55 的“复环面承认交换流”与 p.56 的 Lagrange 论证，是本节最值钱的两句。
\end{itemize}
\end{follow}

\begin{checkbox}
\begin{itemize}
\item 写出 Hitchin 方程的三条，并说明它们分别是哪三个矩映射的零点。
\item 复结构 $I$ 与 $J$ 下，把全纯矩映射设为零再商 $G_\CC$，分别得到什么几何对象？两条定理各叫什么名字？
\item 为什么 $\dim_\CC\MH=(2g-2)\dim G$ 而 $\dim B=(g-1)\dim G$？两者相差的那一半落在哪里？
\item “纤维是复环面”与“纤维是 Lagrange”各用了什么条件？为什么后者同时是对 $\omega_J$ 与 $\omega_K$ 的陈述？
\item 本讲中哪一句只说明“结构已经造好”，哪一句才是“对偶已经做出来”？§4.3 里有没有 T-对偶？
\end{itemize}
\end{checkbox}
